% Options for packages loaded elsewhere
\PassOptionsToPackage{unicode}{hyperref}
\PassOptionsToPackage{hyphens}{url}
\PassOptionsToPackage{space}{xeCJK}
\documentclass[
]{article}
\usepackage{xcolor}
\usepackage[margin=1in]{geometry}
\usepackage{amsmath,amssymb}
\setcounter{secnumdepth}{-\maxdimen} % remove section numbering
\usepackage{iftex}
\ifPDFTeX
  \usepackage[T1]{fontenc}
  \usepackage[utf8]{inputenc}
  \usepackage{textcomp} % provide euro and other symbols
\else % if luatex or xetex
  \usepackage{unicode-math} % this also loads fontspec
  \defaultfontfeatures{Scale=MatchLowercase}
  \defaultfontfeatures[\rmfamily]{Ligatures=TeX,Scale=1}
\fi
\usepackage{lmodern}
\ifPDFTeX\else
  % xetex/luatex font selection
  \ifXeTeX
    \usepackage{xeCJK}
    \setCJKmainfont[]{Noto Sans CJK JP}
  \fi
  \ifLuaTeX
    \usepackage[]{luatexja-fontspec}
    \setmainjfont[]{Noto Sans CJK JP}
  \fi
\fi
% Use upquote if available, for straight quotes in verbatim environments
\IfFileExists{upquote.sty}{\usepackage{upquote}}{}
\IfFileExists{microtype.sty}{% use microtype if available
  \usepackage[]{microtype}
  \UseMicrotypeSet[protrusion]{basicmath} % disable protrusion for tt fonts
}{}
\makeatletter
\@ifundefined{KOMAClassName}{% if non-KOMA class
  \IfFileExists{parskip.sty}{%
    \usepackage{parskip}
  }{% else
    \setlength{\parindent}{0pt}
    \setlength{\parskip}{6pt plus 2pt minus 1pt}}
}{% if KOMA class
  \KOMAoptions{parskip=half}}
\makeatother
\usepackage{color}
\usepackage{fancyvrb}
\newcommand{\VerbBar}{|}
\newcommand{\VERB}{\Verb[commandchars=\\\{\}]}
\DefineVerbatimEnvironment{Highlighting}{Verbatim}{commandchars=\\\{\}}
% Add ',fontsize=\small' for more characters per line
\usepackage{framed}
\definecolor{shadecolor}{RGB}{248,248,248}
\newenvironment{Shaded}{\begin{snugshade}}{\end{snugshade}}
\newcommand{\AlertTok}[1]{\textcolor[rgb]{0.94,0.16,0.16}{#1}}
\newcommand{\AnnotationTok}[1]{\textcolor[rgb]{0.56,0.35,0.01}{\textbf{\textit{#1}}}}
\newcommand{\AttributeTok}[1]{\textcolor[rgb]{0.13,0.29,0.53}{#1}}
\newcommand{\BaseNTok}[1]{\textcolor[rgb]{0.00,0.00,0.81}{#1}}
\newcommand{\BuiltInTok}[1]{#1}
\newcommand{\CharTok}[1]{\textcolor[rgb]{0.31,0.60,0.02}{#1}}
\newcommand{\CommentTok}[1]{\textcolor[rgb]{0.56,0.35,0.01}{\textit{#1}}}
\newcommand{\CommentVarTok}[1]{\textcolor[rgb]{0.56,0.35,0.01}{\textbf{\textit{#1}}}}
\newcommand{\ConstantTok}[1]{\textcolor[rgb]{0.56,0.35,0.01}{#1}}
\newcommand{\ControlFlowTok}[1]{\textcolor[rgb]{0.13,0.29,0.53}{\textbf{#1}}}
\newcommand{\DataTypeTok}[1]{\textcolor[rgb]{0.13,0.29,0.53}{#1}}
\newcommand{\DecValTok}[1]{\textcolor[rgb]{0.00,0.00,0.81}{#1}}
\newcommand{\DocumentationTok}[1]{\textcolor[rgb]{0.56,0.35,0.01}{\textbf{\textit{#1}}}}
\newcommand{\ErrorTok}[1]{\textcolor[rgb]{0.64,0.00,0.00}{\textbf{#1}}}
\newcommand{\ExtensionTok}[1]{#1}
\newcommand{\FloatTok}[1]{\textcolor[rgb]{0.00,0.00,0.81}{#1}}
\newcommand{\FunctionTok}[1]{\textcolor[rgb]{0.13,0.29,0.53}{\textbf{#1}}}
\newcommand{\ImportTok}[1]{#1}
\newcommand{\InformationTok}[1]{\textcolor[rgb]{0.56,0.35,0.01}{\textbf{\textit{#1}}}}
\newcommand{\KeywordTok}[1]{\textcolor[rgb]{0.13,0.29,0.53}{\textbf{#1}}}
\newcommand{\NormalTok}[1]{#1}
\newcommand{\OperatorTok}[1]{\textcolor[rgb]{0.81,0.36,0.00}{\textbf{#1}}}
\newcommand{\OtherTok}[1]{\textcolor[rgb]{0.56,0.35,0.01}{#1}}
\newcommand{\PreprocessorTok}[1]{\textcolor[rgb]{0.56,0.35,0.01}{\textit{#1}}}
\newcommand{\RegionMarkerTok}[1]{#1}
\newcommand{\SpecialCharTok}[1]{\textcolor[rgb]{0.81,0.36,0.00}{\textbf{#1}}}
\newcommand{\SpecialStringTok}[1]{\textcolor[rgb]{0.31,0.60,0.02}{#1}}
\newcommand{\StringTok}[1]{\textcolor[rgb]{0.31,0.60,0.02}{#1}}
\newcommand{\VariableTok}[1]{\textcolor[rgb]{0.00,0.00,0.00}{#1}}
\newcommand{\VerbatimStringTok}[1]{\textcolor[rgb]{0.31,0.60,0.02}{#1}}
\newcommand{\WarningTok}[1]{\textcolor[rgb]{0.56,0.35,0.01}{\textbf{\textit{#1}}}}
\usepackage{longtable,booktabs,array}
\newcounter{none} % for unnumbered tables
\usepackage{calc} % for calculating minipage widths
% Correct order of tables after \paragraph or \subparagraph
\usepackage{etoolbox}
\makeatletter
\patchcmd\longtable{\par}{\if@noskipsec\mbox{}\fi\par}{}{}
\makeatother
% Allow footnotes in longtable head/foot
\IfFileExists{footnotehyper.sty}{\usepackage{footnotehyper}}{\usepackage{footnote}}
\makesavenoteenv{longtable}
\usepackage{graphicx}
\makeatletter
\newsavebox\pandoc@box
\newcommand*\pandocbounded[1]{% scales image to fit in text height/width
  \sbox\pandoc@box{#1}%
  \Gscale@div\@tempa{\textheight}{\dimexpr\ht\pandoc@box+\dp\pandoc@box\relax}%
  \Gscale@div\@tempb{\linewidth}{\wd\pandoc@box}%
  \ifdim\@tempb\p@<\@tempa\p@\let\@tempa\@tempb\fi% select the smaller of both
  \ifdim\@tempa\p@<\p@\scalebox{\@tempa}{\usebox\pandoc@box}%
  \else\usebox{\pandoc@box}%
  \fi%
}
% Set default figure placement to htbp
\def\fps@figure{htbp}
\makeatother
\setlength{\emergencystretch}{3em} % prevent overfull lines
\providecommand{\tightlist}{%
  \setlength{\itemsep}{0pt}\setlength{\parskip}{0pt}}
\usepackage{bookmark}
\IfFileExists{xurl.sty}{\usepackage{xurl}}{} % add URL line breaks if available
\urlstyle{same}
% fallback for those not using the hyperref driver hyperxmp:
\makeatletter
\@ifundefined{xmpquote}{\newcommand{\xmpquote}[1]{#1}}{}
\makeatother
\hypersetup{
  pdftitle={Finding a Great Place to Eat in Kyoto},
  pdfauthor={Agassi Iu},
  hidelinks,
  pdfcreator={LaTeX via pandoc}}

\title{Finding a Great Place to Eat in Kyoto}
\usepackage{etoolbox}
\makeatletter
\providecommand{\subtitle}[1]{% add subtitle to \maketitle
  \apptocmd{\@title}{\par {\large #1 \par}}{}{}
}
\makeatother
\subtitle{Lecture 15 live demo: reference analysis}
\author{Agassi Iu}
\date{November 16, 2026}

\begin{document}
\maketitle

{
\setcounter{tocdepth}{2}
\tableofcontents
}
\section{1. The question}\label{the-question}

\textbf{Where should I eat dinner in Kyoto?}

That is too vague to answer directly, so we break it into three smaller
questions:

\begin{enumerate}
\def\labelenumi{\arabic{enumi}.}
\tightlist
\item
  What are the three most popular cuisines? \emph{(exploratory)}
\item
  What are the three best-rated restaurants at each price point?
  \emph{(descriptive)}
\item
  Which restaurants are rated \textbf{much higher} than their price and
  cuisine would lead us to expect? \emph{(predictive: a model, then a
  look at its misses)}
\end{enumerate}

\section{2. The data}\label{the-data}

The data are restaurant reviews from Tabelog, a Japanese review site
(via Kaggle). Put \texttt{Kyoto\_Restaurant\_Info.csv} in the same
folder as this file.

\begin{Shaded}
\begin{Highlighting}[]
\NormalTok{kyoto\_raw }\OtherTok{\textless{}{-}} \FunctionTok{read\_csv}\NormalTok{(}\StringTok{"Kyoto\_Restaurant\_Info.csv"}\NormalTok{)}
\FunctionTok{glimpse}\NormalTok{(kyoto\_raw)}
\end{Highlighting}
\end{Shaded}

\begin{verbatim}
## Rows: 895
## Columns: 14
## $ ...1           <dbl> 1, 2, 3, 4, 5, 6, 7, 8, 9, 10, 11, 12, 13, 14, 15, 16, ~
## $ Name           <chr> "Orudeidainingurajou", "Steak Frites Gaspard zinzin", "~
## $ JapaneseName   <chr> "オールデイダイニング　ラジョウ", "ステックフリット ガスパール ザンザン", "和馬", "お好み焼き ~
## $ Station        <chr> "Kyoto", "Karasuma", "Sanjo", "Tambaguchi", "Kyoto Shiy~
## $ FirstCategory  <chr> "Buffet style", "Bistro", "Izakaya (Tavern)", "Okonomiy~
## $ SecondCategory <chr> "Cafe", "Steak", "Japanese food (other)", "Izakaya (Tav~
## $ DinnerPrice    <chr> "￥4000～￥4999", "￥3000～￥3999", "￥3000～￥3999", "￥3000～￥39~
## $ LunchPrice     <chr> "￥2000～￥2999", "￥1000～￥1999", NA, NA, "￥1000～￥1999", "￥~
## $ TotalRating    <dbl> 3.39, 3.18, 3.28, 3.14, 3.16, 3.08, 3.04, 3.34, 3.03, 3~
## $ DinnerRating   <dbl> 3.20, 3.06, 3.28, 3.14, 3.16, 3.00, 3.05, 3.34, 3.04, 3~
## $ LunchRating    <dbl> 3.38, 3.33, NA, NA, 3.05, 3.11, 3.00, NA, NA, NA, NA, 3~
## $ ReviewNum      <dbl> 56, 70, 7, 16, 23, 24, 22, 23, 3, 92, 4, 120, 2, 39, 13~
## $ Lat            <dbl> 34.98618, 35.00583, 35.00658, 34.99592, 35.00790, 34.98~
## $ Long           <dbl> 135.7613, 135.7599, 135.7702, 135.7483, 135.7682, 135.7~
\end{verbatim}

Before trusting this data, ask the usual questions: who collected it,
who is in it, and who is missing?

\begin{Shaded}
\begin{Highlighting}[]
\CommentTok{\# Ratings: what is the range?}
\FunctionTok{range}\NormalTok{(kyoto\_raw}\SpecialCharTok{$}\NormalTok{DinnerRating)}
\end{Highlighting}
\end{Shaded}

\begin{verbatim}
## [1] 3.00 4.16
\end{verbatim}

\begin{Shaded}
\begin{Highlighting}[]
\CommentTok{\# How many reviews does a typical restaurant have?}
\FunctionTok{summary}\NormalTok{(kyoto\_raw}\SpecialCharTok{$}\NormalTok{ReviewNum)}
\end{Highlighting}
\end{Shaded}

\begin{verbatim}
##    Min. 1st Qu.  Median    Mean 3rd Qu.    Max. 
##    1.00    7.00   15.00   27.08   30.00 1319.00
\end{verbatim}

Two things to notice:

\begin{itemize}
\tightlist
\item
  The lowest rating is exactly 3.0. This suggests the dataset only
  includes restaurants rated 3.0 or above, so \textbf{poorly rated
  restaurants are missing}. We can only compare good and very good
  places.
\item
  Many restaurants have very few reviews, so their ratings are not
  reliable.
\end{itemize}

\subsection{Cleaning}\label{cleaning}

The prices are text such as \texttt{￥3000～￥3999}. We turn each one
into a number: the \textbf{lower end} of its price band.

\begin{Shaded}
\begin{Highlighting}[]
\NormalTok{kyoto }\OtherTok{\textless{}{-}}\NormalTok{ kyoto\_raw }\SpecialCharTok{|\textgreater{}}
  \FunctionTok{select}\NormalTok{(}\SpecialCharTok{{-}}\DecValTok{1}\NormalTok{) }\SpecialCharTok{|\textgreater{}}                                   \CommentTok{\# drop the unnamed row{-}number column}
  \FunctionTok{mutate}\NormalTok{(}
    \CommentTok{\# lower end of the price band; the band "～￥999" has no lower end, so it is 0}
    \AttributeTok{price\_low =} \FunctionTok{if\_else}\NormalTok{(}
      \FunctionTok{str\_starts}\NormalTok{(}\FunctionTok{str\_trim}\NormalTok{(DinnerPrice), }\StringTok{"～"}\NormalTok{),}
      \DecValTok{0}\NormalTok{,}
      \FunctionTok{as.numeric}\NormalTok{(}\FunctionTok{str\_extract}\NormalTok{(DinnerPrice, }\StringTok{"[0{-}9]+"}\NormalTok{))}
\NormalTok{    ),}
    \CommentTok{\# price bands in order from cheapest to most expensive (for plots)}
    \AttributeTok{price\_band =} \FunctionTok{fct\_reorder}\NormalTok{(DinnerPrice, price\_low),}
    \CommentTok{\# the three price points from the question}
    \AttributeTok{price\_tier =} \FunctionTok{case\_when}\NormalTok{(}
\NormalTok{      price\_low }\SpecialCharTok{\textless{}} \DecValTok{3000}  \SpecialCharTok{\textasciitilde{}} \StringTok{"Cheap"}\NormalTok{,      }\CommentTok{\# 0 to 2,999 yen}
\NormalTok{      price\_low }\SpecialCharTok{\textless{}} \DecValTok{8000}  \SpecialCharTok{\textasciitilde{}} \StringTok{"Middle"}\NormalTok{,     }\CommentTok{\# 3,000 to 7,999 yen}
\NormalTok{      price\_low }\SpecialCharTok{\textless{}} \DecValTok{15000} \SpecialCharTok{\textasciitilde{}} \StringTok{"Expensive"}\NormalTok{,  }\CommentTok{\# 8,000 to 14,999 yen}
      \ConstantTok{TRUE}              \SpecialCharTok{\textasciitilde{}} \ConstantTok{NA\_character\_} \CommentTok{\# 15,000 yen and above: not in our tiers}
\NormalTok{    ),}
    \AttributeTok{price\_tier =} \FunctionTok{factor}\NormalTok{(price\_tier, }\AttributeTok{levels =} \FunctionTok{c}\NormalTok{(}\StringTok{"Cheap"}\NormalTok{, }\StringTok{"Middle"}\NormalTok{, }\StringTok{"Expensive"}\NormalTok{))}
\NormalTok{  )}

\NormalTok{kyoto }\SpecialCharTok{|\textgreater{}}
  \FunctionTok{count}\NormalTok{(price\_tier) }\SpecialCharTok{|\textgreater{}}
  \FunctionTok{mutate}\NormalTok{(}\AttributeTok{note =} \FunctionTok{if\_else}\NormalTok{(}\FunctionTok{is.na}\NormalTok{(price\_tier), }\StringTok{"15,000 yen and above (left out)"}\NormalTok{, }\StringTok{""}\NormalTok{))}
\end{Highlighting}
\end{Shaded}

{\def\LTcaptype{none} % do not increment counter
\begin{longtable}[]{@{}lrl@{}}
\toprule\noalign{}
price\_tier & n & note \\
\midrule\noalign{}
\endhead
\bottomrule\noalign{}
\endlastfoot
Cheap & 206 & \\
Middle & 596 & \\
Expensive & 79 & \\
NA & 14 & 15,000 yen and above (left out) \\
\end{longtable}
}

\subsection{Keep only restaurants with enough
reviews}\label{keep-only-restaurants-with-enough-reviews}

A rating based on one or two reviews tells us very little. For every
question below we keep restaurants with \textbf{at least 10 reviews}.

\begin{Shaded}
\begin{Highlighting}[]
\NormalTok{rated }\OtherTok{\textless{}{-}}\NormalTok{ kyoto }\SpecialCharTok{|\textgreater{}} \FunctionTok{filter}\NormalTok{(ReviewNum }\SpecialCharTok{\textgreater{}=} \DecValTok{10}\NormalTok{)}

\FunctionTok{c}\NormalTok{(}\AttributeTok{before =} \FunctionTok{nrow}\NormalTok{(kyoto), }\AttributeTok{after =} \FunctionTok{nrow}\NormalTok{(rated))}
\end{Highlighting}
\end{Shaded}

\begin{verbatim}
## before  after 
##    895    602
\end{verbatim}

\section{3. Question 1: the three most popular
cuisines}\label{question-1-the-three-most-popular-cuisines}

``Popular'' can mean different things. We look at two measures and see
whether they agree.

\begin{itemize}
\tightlist
\item
  \textbf{Number of restaurants} in the category
\item
  \textbf{Total number of reviews} across the category (a rough measure
  of how many diners it attracts)
\end{itemize}

\begin{Shaded}
\begin{Highlighting}[]
\NormalTok{popular }\OtherTok{\textless{}{-}}\NormalTok{ rated }\SpecialCharTok{|\textgreater{}}
  \FunctionTok{group\_by}\NormalTok{(FirstCategory) }\SpecialCharTok{|\textgreater{}}
  \FunctionTok{summarise}\NormalTok{(}
    \AttributeTok{restaurants   =} \FunctionTok{n}\NormalTok{(),}
    \AttributeTok{total\_reviews =} \FunctionTok{sum}\NormalTok{(ReviewNum),}
    \AttributeTok{.groups =} \StringTok{"drop"}
\NormalTok{  )}

\NormalTok{popular }\SpecialCharTok{|\textgreater{}} \FunctionTok{arrange}\NormalTok{(}\FunctionTok{desc}\NormalTok{(total\_reviews)) }\SpecialCharTok{|\textgreater{}} \FunctionTok{head}\NormalTok{(}\DecValTok{10}\NormalTok{)}
\end{Highlighting}
\end{Shaded}

{\def\LTcaptype{none} % do not increment counter
\begin{longtable}[]{@{}lrr@{}}
\toprule\noalign{}
FirstCategory & restaurants & total\_reviews \\
\midrule\noalign{}
\endhead
\bottomrule\noalign{}
\endlastfoot
Izakaya (Tavern) & 121 & 2945 \\
Ramen & 12 & 2569 \\
Italian & 54 & 2444 \\
Yakiniku (BBQ Beef) & 54 & 1199 \\
Kyoto Cuisine & 33 & 1059 \\
Kaiseki (Traditional Japanese) & 23 & 894 \\
Bar & 24 & 630 \\
Cafe & 16 & 583 \\
Buffet style & 12 & 581 \\
Tempura & 8 & 534 \\
\end{longtable}
}

\begin{Shaded}
\begin{Highlighting}[]
\NormalTok{top3\_cuisines }\OtherTok{\textless{}{-}}\NormalTok{ popular }\SpecialCharTok{|\textgreater{}}
  \FunctionTok{arrange}\NormalTok{(}\FunctionTok{desc}\NormalTok{(total\_reviews)) }\SpecialCharTok{|\textgreater{}}
  \FunctionTok{slice\_head}\NormalTok{(}\AttributeTok{n =} \DecValTok{3}\NormalTok{)}

\NormalTok{top3\_cuisines}
\end{Highlighting}
\end{Shaded}

{\def\LTcaptype{none} % do not increment counter
\begin{longtable}[]{@{}lrr@{}}
\toprule\noalign{}
FirstCategory & restaurants & total\_reviews \\
\midrule\noalign{}
\endhead
\bottomrule\noalign{}
\endlastfoot
Izakaya (Tavern) & 121 & 2945 \\
Ramen & 12 & 2569 \\
Italian & 54 & 2444 \\
\end{longtable}
}

\begin{Shaded}
\begin{Highlighting}[]
\NormalTok{popular }\SpecialCharTok{|\textgreater{}}
  \FunctionTok{slice\_max}\NormalTok{(total\_reviews, }\AttributeTok{n =} \DecValTok{10}\NormalTok{, }\AttributeTok{with\_ties =} \ConstantTok{FALSE}\NormalTok{) }\SpecialCharTok{|\textgreater{}}
  \FunctionTok{mutate}\NormalTok{(}\AttributeTok{top3 =}\NormalTok{ FirstCategory }\SpecialCharTok{\%in\%}\NormalTok{ top3\_cuisines}\SpecialCharTok{$}\NormalTok{FirstCategory) }\SpecialCharTok{|\textgreater{}}
  \FunctionTok{ggplot}\NormalTok{(}\FunctionTok{aes}\NormalTok{(}\AttributeTok{x =}\NormalTok{ total\_reviews, }\AttributeTok{y =} \FunctionTok{fct\_reorder}\NormalTok{(FirstCategory, total\_reviews), }\AttributeTok{fill =}\NormalTok{ top3)) }\SpecialCharTok{+}
  \FunctionTok{geom\_col}\NormalTok{(}\AttributeTok{show.legend =} \ConstantTok{FALSE}\NormalTok{) }\SpecialCharTok{+}
  \FunctionTok{scale\_fill\_manual}\NormalTok{(}\AttributeTok{values =} \FunctionTok{c}\NormalTok{(}\StringTok{"grey70"}\NormalTok{, }\StringTok{"steelblue"}\NormalTok{)) }\SpecialCharTok{+}
  \FunctionTok{labs}\NormalTok{(}\AttributeTok{x =} \StringTok{"Total reviews"}\NormalTok{, }\AttributeTok{y =} \ConstantTok{NULL}\NormalTok{,}
       \AttributeTok{title =} \StringTok{"Cuisines diners write about most"}\NormalTok{)}
\end{Highlighting}
\end{Shaded}

\begin{center}\includegraphics{Mini-project-walkthrough_files/figure-latex/q1-plot-1} \end{center}

If we count restaurants instead of reviews, does the top three change?

\begin{Shaded}
\begin{Highlighting}[]
\NormalTok{popular }\SpecialCharTok{|\textgreater{}} \FunctionTok{arrange}\NormalTok{(}\FunctionTok{desc}\NormalTok{(restaurants)) }\SpecialCharTok{|\textgreater{}} \FunctionTok{head}\NormalTok{(}\DecValTok{5}\NormalTok{)}
\end{Highlighting}
\end{Shaded}

{\def\LTcaptype{none} % do not increment counter
\begin{longtable}[]{@{}lrr@{}}
\toprule\noalign{}
FirstCategory & restaurants & total\_reviews \\
\midrule\noalign{}
\endhead
\bottomrule\noalign{}
\endlastfoot
Izakaya (Tavern) & 121 & 2945 \\
Italian & 54 & 2444 \\
Yakiniku (BBQ Beef) & 54 & 1199 \\
Kyoto Cuisine & 33 & 1059 \\
Bar & 24 & 630 \\
\end{longtable}
}

\begin{quote}
\textbf{Note:} The answer depends on how we \textbf{define} popular.
Always say which definition you used.
\end{quote}

\section{4. Question 2: the top three restaurants at each price
point}\label{question-2-the-top-three-restaurants-at-each-price-point}

\begin{Shaded}
\begin{Highlighting}[]
\NormalTok{top\_by\_tier }\OtherTok{\textless{}{-}}\NormalTok{ rated }\SpecialCharTok{|\textgreater{}}
  \FunctionTok{filter}\NormalTok{(}\SpecialCharTok{!}\FunctionTok{is.na}\NormalTok{(price\_tier)) }\SpecialCharTok{|\textgreater{}}
  \FunctionTok{arrange}\NormalTok{(}\FunctionTok{desc}\NormalTok{(DinnerRating), }\FunctionTok{desc}\NormalTok{(ReviewNum)) }\SpecialCharTok{|\textgreater{}}   \CommentTok{\# best rating first; more reviews breaks ties}
  \FunctionTok{group\_by}\NormalTok{(price\_tier) }\SpecialCharTok{|\textgreater{}}
  \FunctionTok{slice\_head}\NormalTok{(}\AttributeTok{n =} \DecValTok{3}\NormalTok{) }\SpecialCharTok{|\textgreater{}}
  \FunctionTok{ungroup}\NormalTok{() }\SpecialCharTok{|\textgreater{}}
  \FunctionTok{arrange}\NormalTok{(price\_tier, }\FunctionTok{desc}\NormalTok{(DinnerRating)) }\SpecialCharTok{|\textgreater{}}
  \FunctionTok{select}\NormalTok{(price\_tier, Name, JapaneseName, FirstCategory,}
\NormalTok{         DinnerPrice, DinnerRating, ReviewNum)}

\NormalTok{top\_by\_tier}
\end{Highlighting}
\end{Shaded}

{\def\LTcaptype{none} % do not increment counter
\begin{longtable}[]{@{}
  >{\raggedright\arraybackslash}p{(\linewidth - 12\tabcolsep) * \real{0.0880}}
  >{\raggedright\arraybackslash}p{(\linewidth - 12\tabcolsep) * \real{0.2160}}
  >{\raggedright\arraybackslash}p{(\linewidth - 12\tabcolsep) * \real{0.2240}}
  >{\raggedright\arraybackslash}p{(\linewidth - 12\tabcolsep) * \real{0.1520}}
  >{\raggedright\arraybackslash}p{(\linewidth - 12\tabcolsep) * \real{0.1360}}
  >{\raggedleft\arraybackslash}p{(\linewidth - 12\tabcolsep) * \real{0.1040}}
  >{\raggedleft\arraybackslash}p{(\linewidth - 12\tabcolsep) * \real{0.0800}}@{}}
\toprule\noalign{}
\begin{minipage}[b]{\linewidth}\raggedright
price\_tier
\end{minipage} & \begin{minipage}[b]{\linewidth}\raggedright
Name
\end{minipage} & \begin{minipage}[b]{\linewidth}\raggedright
JapaneseName
\end{minipage} & \begin{minipage}[b]{\linewidth}\raggedright
FirstCategory
\end{minipage} & \begin{minipage}[b]{\linewidth}\raggedright
DinnerPrice
\end{minipage} & \begin{minipage}[b]{\linewidth}\raggedleft
DinnerRating
\end{minipage} & \begin{minipage}[b]{\linewidth}\raggedleft
ReviewNum
\end{minipage} \\
\midrule\noalign{}
\endhead
\bottomrule\noalign{}
\endlastfoot
Cheap & Giwontokiwa & ぎをん 常盤 & Soba & ￥1000～￥1999 & 3.60 &
106 \\
Cheap & Honkedaiichiasahi & 本家 第一旭 たかばし本店 & Ramen & ～￥999 &
3.59 & 1319 \\
Cheap & chuukasobatakayasu & 中華そば 高安 & Ramen & ～￥999 & 3.58 &
713 \\
Middle & Ryuunohige & 龍のひげ & Kyoto Cuisine & ￥5000～￥5999 & 3.74 &
70 \\
Middle & Zasoudouhigashiyamakyouto & THE SODOH HIGASHIYAMA KYOTO &
Italian & ￥6000～￥7999 & 3.61 & 272 \\
Middle & Takasegawa & TAKASEGAWA & Sumibiyaki & ￥5000～￥5999 & 3.60 &
43 \\
Expensive & Orudohonkonresutorankyouto & 老香港酒家京都 & Chinese &
￥10000～￥14999 & 4.04 & 148 \\
Expensive & ORTO & オルト & Innovative cuisine & ￥8000～￥9999 & 3.97 &
53 \\
Expensive & Iharada & いはら田 & Kyoto Cuisine & ￥10000～￥14999 & 3.91
& 52 \\
\end{longtable}
}

\begin{Shaded}
\begin{Highlighting}[]
\FunctionTok{ggplot}\NormalTok{(top\_by\_tier,}
       \FunctionTok{aes}\NormalTok{(}\AttributeTok{x =}\NormalTok{ DinnerRating, }\AttributeTok{y =} \FunctionTok{fct\_reorder}\NormalTok{(Name, DinnerRating), }\AttributeTok{color =}\NormalTok{ FirstCategory)) }\SpecialCharTok{+}
  \FunctionTok{geom\_point}\NormalTok{(}\AttributeTok{size =} \DecValTok{4}\NormalTok{) }\SpecialCharTok{+}
  \FunctionTok{facet\_wrap}\NormalTok{(}\SpecialCharTok{\textasciitilde{}}\NormalTok{ price\_tier, }\AttributeTok{ncol =} \DecValTok{1}\NormalTok{, }\AttributeTok{scales =} \StringTok{"free\_y"}\NormalTok{) }\SpecialCharTok{+}
  \FunctionTok{labs}\NormalTok{(}\AttributeTok{x =} \StringTok{"Dinner rating"}\NormalTok{, }\AttributeTok{y =} \ConstantTok{NULL}\NormalTok{, }\AttributeTok{color =} \StringTok{"Cuisine"}\NormalTok{,}
       \AttributeTok{title =} \StringTok{"Top three restaurants at each price point"}\NormalTok{)}
\end{Highlighting}
\end{Shaded}

\begin{center}\includegraphics{Mini-project-walkthrough_files/figure-latex/q2-plot-1} \end{center}

\section{5. Question 3: which places beat their
price?}\label{question-3-which-places-beat-their-price}

The top-three lists favour expensive restaurants, because pricier places
tend to be rated higher. A more useful question for a diner is:
\textbf{which restaurants are rated higher than we would expect for
their price and cuisine?}

\subsection{Step 1: build a simple model of the typical
rating}\label{step-1-build-a-simple-model-of-the-typical-rating}

We predict \texttt{DinnerRating} from the price and the cuisine.
Cuisines with fewer than 15 restaurants are grouped as ``Other'' so the
model has enough data for each one.

\begin{Shaded}
\begin{Highlighting}[]
\NormalTok{model\_data }\OtherTok{\textless{}{-}}\NormalTok{ rated }\SpecialCharTok{|\textgreater{}}
  \FunctionTok{mutate}\NormalTok{(}
    \AttributeTok{price\_k =}\NormalTok{ price\_low }\SpecialCharTok{/} \DecValTok{1000}\NormalTok{,                      }\CommentTok{\# price in thousands of yen}
    \AttributeTok{cuisine =} \FunctionTok{fct\_lump\_min}\NormalTok{(FirstCategory, }\AttributeTok{min =} \DecValTok{15}\NormalTok{)  }\CommentTok{\# small cuisines become "Other"}
\NormalTok{  )}

\NormalTok{rating\_fit }\OtherTok{\textless{}{-}} \FunctionTok{lm}\NormalTok{(DinnerRating }\SpecialCharTok{\textasciitilde{}}\NormalTok{ price\_k }\SpecialCharTok{+}\NormalTok{ cuisine, }\AttributeTok{data =}\NormalTok{ model\_data)}

\NormalTok{broom}\SpecialCharTok{::}\FunctionTok{tidy}\NormalTok{(rating\_fit)}
\end{Highlighting}
\end{Shaded}

{\def\LTcaptype{none} % do not increment counter
\begin{longtable}[]{@{}
  >{\raggedright\arraybackslash}p{(\linewidth - 8\tabcolsep) * \real{0.4750}}
  >{\raggedleft\arraybackslash}p{(\linewidth - 8\tabcolsep) * \real{0.1375}}
  >{\raggedleft\arraybackslash}p{(\linewidth - 8\tabcolsep) * \real{0.1250}}
  >{\raggedleft\arraybackslash}p{(\linewidth - 8\tabcolsep) * \real{0.1375}}
  >{\raggedleft\arraybackslash}p{(\linewidth - 8\tabcolsep) * \real{0.1250}}@{}}
\toprule\noalign{}
\begin{minipage}[b]{\linewidth}\raggedright
term
\end{minipage} & \begin{minipage}[b]{\linewidth}\raggedleft
estimate
\end{minipage} & \begin{minipage}[b]{\linewidth}\raggedleft
std.error
\end{minipage} & \begin{minipage}[b]{\linewidth}\raggedleft
statistic
\end{minipage} & \begin{minipage}[b]{\linewidth}\raggedleft
p.value
\end{minipage} \\
\midrule\noalign{}
\endhead
\bottomrule\noalign{}
\endlastfoot
(Intercept) & 3.1592615 & 0.0403123 & 78.3697619 & 0.0000000 \\
price\_k & 0.0286363 & 0.0028435 & 10.0709319 & 0.0000000 \\
cuisineCafe & -0.1322160 & 0.0623430 & -2.1207833 & 0.0343559 \\
cuisineDining bar & -0.0438068 & 0.0634251 & -0.6906864 & 0.4900338 \\
cuisineFrench & -0.1156515 & 0.0625090 & -1.8501589 & 0.0647896 \\
cuisineItalian & -0.0542445 & 0.0474509 & -1.1431700 & 0.2534309 \\
cuisineIzakaya (Tavern) & -0.0907753 & 0.0430440 & -2.1088931 &
0.0353739 \\
cuisineKaiseki (Traditional Japanese) & -0.0390042 & 0.0582337 &
-0.6697868 & 0.5032553 \\
cuisineKyoto Cuisine & -0.0199637 & 0.0527338 & -0.3785744 &
0.7051399 \\
cuisineYakiniku (BBQ Beef) & -0.0630829 & 0.0473028 & -1.3335968 &
0.1828498 \\
cuisineOther & -0.0317715 & 0.0412273 & -0.7706425 & 0.4412268 \\
\end{longtable}
}

\begin{quote}
\textbf{Note:} The \texttt{price\_k} row is the change in the predicted
rating for every extra 1,000 yen on the price band, with cuisine held
fixed. The cuisine rows compare each cuisine with the reference cuisine
(the first level).
\end{quote}

The model explains only part of the variation in ratings:

\begin{Shaded}
\begin{Highlighting}[]
\FunctionTok{summary}\NormalTok{(rating\_fit)}\SpecialCharTok{$}\NormalTok{r.squared}
\end{Highlighting}
\end{Shaded}

\begin{verbatim}
## [1] 0.2040006
\end{verbatim}

This is expected. Price and cuisine alone cannot predict how good a
restaurant is, and that leftover is exactly what we want to look at.

\subsection{Step 2: look at the misses}\label{step-2-look-at-the-misses}

For each restaurant, the \textbf{gap} is its actual rating minus the
rating the model expected. A large positive gap means it is rated much
higher than its price and cuisine suggest.

\begin{Shaded}
\begin{Highlighting}[]
\NormalTok{gems }\OtherTok{\textless{}{-}}\NormalTok{ model\_data }\SpecialCharTok{|\textgreater{}}
  \FunctionTok{mutate}\NormalTok{(}
    \AttributeTok{expected =} \FunctionTok{fitted}\NormalTok{(rating\_fit),}
    \AttributeTok{gap      =}\NormalTok{ DinnerRating }\SpecialCharTok{{-}}\NormalTok{ expected}
\NormalTok{  )}

\NormalTok{gems }\SpecialCharTok{|\textgreater{}}
  \FunctionTok{arrange}\NormalTok{(}\FunctionTok{desc}\NormalTok{(gap)) }\SpecialCharTok{|\textgreater{}}
  \FunctionTok{select}\NormalTok{(Name, JapaneseName, FirstCategory, DinnerPrice,}
\NormalTok{         DinnerRating, expected, gap, ReviewNum) }\SpecialCharTok{|\textgreater{}}
  \FunctionTok{mutate}\NormalTok{(}\FunctionTok{across}\NormalTok{(}\FunctionTok{c}\NormalTok{(expected, gap), \textbackslash{}(x) }\FunctionTok{round}\NormalTok{(x, }\DecValTok{2}\NormalTok{))) }\SpecialCharTok{|\textgreater{}}
  \FunctionTok{head}\NormalTok{(}\DecValTok{10}\NormalTok{)}
\end{Highlighting}
\end{Shaded}

{\def\LTcaptype{none} % do not increment counter
\begin{longtable}[]{@{}
  >{\raggedright\arraybackslash}p{(\linewidth - 14\tabcolsep) * \real{0.1971}}
  >{\raggedright\arraybackslash}p{(\linewidth - 14\tabcolsep) * \real{0.1825}}
  >{\raggedright\arraybackslash}p{(\linewidth - 14\tabcolsep) * \real{0.2263}}
  >{\raggedright\arraybackslash}p{(\linewidth - 14\tabcolsep) * \real{0.1241}}
  >{\raggedleft\arraybackslash}p{(\linewidth - 14\tabcolsep) * \real{0.0949}}
  >{\raggedleft\arraybackslash}p{(\linewidth - 14\tabcolsep) * \real{0.0657}}
  >{\raggedleft\arraybackslash}p{(\linewidth - 14\tabcolsep) * \real{0.0365}}
  >{\raggedleft\arraybackslash}p{(\linewidth - 14\tabcolsep) * \real{0.0730}}@{}}
\toprule\noalign{}
\begin{minipage}[b]{\linewidth}\raggedright
Name
\end{minipage} & \begin{minipage}[b]{\linewidth}\raggedright
JapaneseName
\end{minipage} & \begin{minipage}[b]{\linewidth}\raggedright
FirstCategory
\end{minipage} & \begin{minipage}[b]{\linewidth}\raggedright
DinnerPrice
\end{minipage} & \begin{minipage}[b]{\linewidth}\raggedleft
DinnerRating
\end{minipage} & \begin{minipage}[b]{\linewidth}\raggedleft
expected
\end{minipage} & \begin{minipage}[b]{\linewidth}\raggedleft
gap
\end{minipage} & \begin{minipage}[b]{\linewidth}\raggedleft
ReviewNum
\end{minipage} \\
\midrule\noalign{}
\endhead
\bottomrule\noalign{}
\endlastfoot
Orudohonkonresutorankyouto & 老香港酒家京都 & Chinese & ￥10000～￥14999
& 4.04 & 3.41 & 0.63 & 148 \\
ORTO & オルト & Innovative cuisine & ￥8000～￥9999 & 3.97 & 3.36 & 0.61
& 53 \\
Iharada & いはら田 & Kyoto Cuisine & ￥10000～￥14999 & 3.91 & 3.43 &
0.48 & 52 \\
Kichisen & 京懐石 吉泉 & Kaiseki (Traditional Japanese) &
￥20000～￥29999 & 4.16 & 3.69 & 0.47 & 79 \\
Honkedaiichiasahi & 本家 第一旭 たかばし本店 & Ramen & ～￥999 & 3.59 &
3.13 & 0.46 & 1319 \\
Ryuunohige & 龍のひげ & Kyoto Cuisine & ￥5000～￥5999 & 3.74 & 3.28 &
0.46 & 70 \\
chuukasobatakayasu & 中華そば 高安 & Ramen & ～￥999 & 3.58 & 3.13 &
0.45 & 713 \\
Giwontokiwa & ぎをん 常盤 & Soba & ￥1000～￥1999 & 3.60 & 3.16 & 0.44 &
106 \\
gionni-yongo & 祗園245 & Italian & ￥10000～￥14999 & 3.83 & 3.39 & 0.44
& 39 \\
Yumeya & 夢屋 & Okonomiyaki & ￥1000～￥1999 & 3.55 & 3.16 & 0.39 &
73 \\
\end{longtable}
}

\begin{Shaded}
\begin{Highlighting}[]
\FunctionTok{ggplot}\NormalTok{(gems, }\FunctionTok{aes}\NormalTok{(}\AttributeTok{x =}\NormalTok{ expected, }\AttributeTok{y =}\NormalTok{ DinnerRating)) }\SpecialCharTok{+}
  \FunctionTok{geom\_point}\NormalTok{(}\AttributeTok{alpha =} \FloatTok{0.4}\NormalTok{) }\SpecialCharTok{+}
  \FunctionTok{geom\_abline}\NormalTok{(}\AttributeTok{linetype =} \StringTok{"dashed"}\NormalTok{) }\SpecialCharTok{+}
  \FunctionTok{geom\_point}\NormalTok{(}\AttributeTok{data =} \FunctionTok{slice\_max}\NormalTok{(gems, gap, }\AttributeTok{n =} \DecValTok{5}\NormalTok{, }\AttributeTok{with\_ties =} \ConstantTok{FALSE}\NormalTok{),}
             \AttributeTok{color =} \StringTok{"firebrick"}\NormalTok{, }\AttributeTok{size =} \DecValTok{3}\NormalTok{) }\SpecialCharTok{+}
  \FunctionTok{geom\_text}\NormalTok{(}\AttributeTok{data =} \FunctionTok{slice\_max}\NormalTok{(gems, gap, }\AttributeTok{n =} \DecValTok{5}\NormalTok{, }\AttributeTok{with\_ties =} \ConstantTok{FALSE}\NormalTok{),}
            \FunctionTok{aes}\NormalTok{(}\AttributeTok{label =}\NormalTok{ Name), }\AttributeTok{hjust =} \FloatTok{1.05}\NormalTok{, }\AttributeTok{size =} \FloatTok{3.2}\NormalTok{, }\AttributeTok{check\_overlap =} \ConstantTok{TRUE}\NormalTok{) }\SpecialCharTok{+}
  \FunctionTok{labs}\NormalTok{(}\AttributeTok{x =} \StringTok{"Rating expected from price and cuisine"}\NormalTok{,}
       \AttributeTok{y =} \StringTok{"Actual dinner rating"}\NormalTok{,}
       \AttributeTok{title =} \StringTok{"Points above the dashed line beat expectations"}\NormalTok{)}
\end{Highlighting}
\end{Shaded}

\begin{center}\includegraphics{Mini-project-walkthrough_files/figure-latex/q3-plot-1} \end{center}

\subsection{Step 3: hidden gems on a
budget}\label{step-3-hidden-gems-on-a-budget}

\begin{Shaded}
\begin{Highlighting}[]
\NormalTok{gems }\SpecialCharTok{|\textgreater{}}
  \FunctionTok{filter}\NormalTok{(price\_tier }\SpecialCharTok{==} \StringTok{"Cheap"}\NormalTok{) }\SpecialCharTok{|\textgreater{}}
  \FunctionTok{arrange}\NormalTok{(}\FunctionTok{desc}\NormalTok{(gap)) }\SpecialCharTok{|\textgreater{}}
  \FunctionTok{select}\NormalTok{(Name, JapaneseName, FirstCategory, DinnerPrice, DinnerRating, gap, ReviewNum) }\SpecialCharTok{|\textgreater{}}
  \FunctionTok{mutate}\NormalTok{(}\AttributeTok{gap =} \FunctionTok{round}\NormalTok{(gap, }\DecValTok{2}\NormalTok{)) }\SpecialCharTok{|\textgreater{}}
  \FunctionTok{head}\NormalTok{(}\DecValTok{5}\NormalTok{)}
\end{Highlighting}
\end{Shaded}

{\def\LTcaptype{none} % do not increment counter
\begin{longtable}[]{@{}
  >{\raggedright\arraybackslash}p{(\linewidth - 12\tabcolsep) * \real{0.1881}}
  >{\raggedright\arraybackslash}p{(\linewidth - 12\tabcolsep) * \real{0.2475}}
  >{\raggedright\arraybackslash}p{(\linewidth - 12\tabcolsep) * \real{0.1386}}
  >{\raggedright\arraybackslash}p{(\linewidth - 12\tabcolsep) * \real{0.1485}}
  >{\raggedleft\arraybackslash}p{(\linewidth - 12\tabcolsep) * \real{0.1287}}
  >{\raggedleft\arraybackslash}p{(\linewidth - 12\tabcolsep) * \real{0.0495}}
  >{\raggedleft\arraybackslash}p{(\linewidth - 12\tabcolsep) * \real{0.0990}}@{}}
\toprule\noalign{}
\begin{minipage}[b]{\linewidth}\raggedright
Name
\end{minipage} & \begin{minipage}[b]{\linewidth}\raggedright
JapaneseName
\end{minipage} & \begin{minipage}[b]{\linewidth}\raggedright
FirstCategory
\end{minipage} & \begin{minipage}[b]{\linewidth}\raggedright
DinnerPrice
\end{minipage} & \begin{minipage}[b]{\linewidth}\raggedleft
DinnerRating
\end{minipage} & \begin{minipage}[b]{\linewidth}\raggedleft
gap
\end{minipage} & \begin{minipage}[b]{\linewidth}\raggedleft
ReviewNum
\end{minipage} \\
\midrule\noalign{}
\endhead
\bottomrule\noalign{}
\endlastfoot
Honkedaiichiasahi & 本家 第一旭 たかばし本店 & Ramen & ～￥999 & 3.59 &
0.46 & 1319 \\
chuukasobatakayasu & 中華そば 高安 & Ramen & ～￥999 & 3.58 & 0.45 &
713 \\
Giwontokiwa & ぎをん 常盤 & Soba & ￥1000～￥1999 & 3.60 & 0.44 & 106 \\
Yumeya & 夢屋 & Okonomiyaki & ￥1000～￥1999 & 3.55 & 0.39 & 73 \\
Suiba & すいば 蛸薬師室町店 & Stand Bar & ￥1000～￥1999 & 3.53 & 0.37 &
29 \\
\end{longtable}
}

\section{6. What we found}\label{what-we-found}

\begin{itemize}
\tightlist
\item
  The three most popular cuisines, measured by total reviews, are
  \textbf{Izakaya (Tavern), Ramen, Italian}.
\item
  At each price point the best-rated places are in the table in Question
  2.
\item
  The restaurants that beat their price and cuisine by the most are in
  the table in Question 3. They are good candidates for ``where should I
  eat?''.
\end{itemize}

\subsection{What this analysis can and cannot
say}\label{what-this-analysis-can-and-cannot-say}

\begin{itemize}
\tightlist
\item
  \textbf{Missing bad restaurants.} Ratings start at 3.0, so we only see
  places that already rate well.
\item
  \textbf{Popularity is not quality.} Many reviews means many diners,
  not necessarily better food.
\item
  \textbf{Associations only.} The model describes how price, cuisine and
  rating go together. It does not say that paying more \emph{causes} a
  better rating.
\item
  \textbf{The model is simple.} It uses two variables and leaves a lot
  of variation unexplained, which is why the gaps are interesting.
\item
  \textbf{A decision for you.} The tiers use the \emph{lower end} of
  each price band, and restaurants priced at 15,000 yen and above are
  left out.
\end{itemize}

\end{document}
